\documentclass[10pt]{article}
\usepackage[T1]{fontenc}
\usepackage[utf8]{inputenc}
\usepackage[french]{babel}
\usepackage{amsmath,amssymb,mathrsfs,mathpazo}
\usepackage{geometry}
\geometry{a4paper,margin=2cm}
\usepackage{graphicx,array,tabularx,tabularray,colortbl,multirow,multicol,booktabs,enumitem,pifont,ulem,fancybox,diagbox,cancel}
\usepackage{pstricks,pstricks-add,pst-plot,pst-tree,pst-node,pst-text,pst-func,pst-math,pst-bspline,pst-3dplot,pst-eucl,pst-solides3d}
\let\pscalculate\undefined
\usepackage{pst-bezier}
\usepackage{tikz,pgfplots,tkz-tab}
\usetikzlibrary{arrows,arrows.meta,patterns,calc,positioning,trees,decorations.markings,shapes.geometric,shapes.misc,shapes.arrows,intersections,decorations.pathmorphing,decorations.pathreplacing,backgrounds,shadows,babel}
\pgfplotsset{compat=1.18}
\usepgfplotslibrary{fillbetween}
\usepackage[np]{numprint}
\usepackage{listings,makecell,pas-tableur,twoopt}
\let\barre\undefined
\usepackage[pstricks]{tablvar}
\definecolor{darkgreen}{RGB}{0,100,0}
\providecommand{\red}{\color{red}}
\providecommand{\blue}{\color{blue}}
\providecommand{\green}{\color{green}}
\providecommand{\black}{\color{black}}
\providecommand{\white}{\color{white}}
\providecommand{\gray}{\color{gray}}
\providecommand{\cyan}{\color{cyan}}
\providecommand{\magenta}{\color{magenta}}
\providecommand{\yellow}{\color{yellow}}
\providecommand{\violet}{\color{violet}}
\providecommand{\orange}{\color{orange}}
\providecommand{\euro}{\texteuro}
\providecommand{\EUR}{\texteuro}
\providecommand{\ieme}{\textsuperscript{e}}
\providecommand{\degres}{\ensuremath{^\circ}}
\providecommand{\degree}{\ensuremath{^\circ}}
\providecommand{\vv}[1]{\overrightarrow{#1}}
\providecommand{\Coord}[1]{#1}
\providecommand{\decosix}{}
\providecommand{\starredbullet}{\ensuremath{\bullet}}
\pagestyle{empty}
\setlength{\parindent}{0pt}
\usepackage{adjustbox}\DeclareUnicodeCharacter{2236}{:}\DeclareUnicodeCharacter{2190}{\ensuremath{\leftarrow}}
\providecommand{\R}{\mathbb{R}}
\providecommand{\N}{\mathbb{N}}
\providecommand{\Z}{\mathbb{Z}}
\providecommand{\Q}{\mathbb{Q}}
\providecommand{\D}{\mathbb{D}}
\providecommand{\C}{\mathbb{C}}
\providecommand{\e}{{\mathrm{e}}}
\providecommand{\vect}[1]{\overrightarrow{#1}}
\providecommand{\vectt}[1]{\overrightarrow{\mathrm{#1}}}
\providecommand{\vvt}[1]{\overrightarrow{\mathrm{#1}}}
\providecommand{\np}[2][]{{#2~#1}}
\providecommand{\vv}[1]{{\overrightarrow{#1}}}
\renewcommand{\d}{{\mathrm{d}}}
\providecommand{\degres}{{^\circ}}
\providecommand{\degre}{{^\circ}}
\providecommand{\blue}{\color{blue}}
\providecommand{\red}{\color{red}}
\providecommand{\enskip}{\quad}
\providecommand{\llbracket}{\lbrack\!\lbrack}
\providecommand{\rrbracket}{\rbrack\!\rbrack}
\providecommand{\euro}{{\text{€}}}
\providecommand{\no}{{\text{n°}}}
\providecommand{\interval}[2]{{[#1;#2]}}
\providecommand{\intervFO}[2]{{[#1;#2[}}
\providecommand{\dotfill}{\ldots}

\providecommand{\ds}{\displaystyle}
\providecommand{\textding}[1]{#1}
\providecommand{\Oij}{\(\mathrm{O};\vec{\imath},\vec{\jmath}\)}
\providecommand{\Oijk}{\(\mathrm{O};\vec{\imath},\vec{\jmath},\vec{k}\)}
\providecommand{\Ouv}{\(\mathrm{O};\vec{u},\vec{v}\)}
\providecommand{\intervFF}[2]{\ensuremath{[#1;#2]}}
\renewcommand{\euro}{\ifmmode\text{\texteuro}\else\texteuro\fi}
\DeclareUnicodeCharacter{00B0}{\ifmmode^{\circ}\else\textdegree\fi}
\DeclareUnicodeCharacter{20AC}{\ifmmode\text{\texteuro}\else\texteuro\fi}
% Énoncés : Éducation nationale. Transcription et corrigés : APMEP, Denis Vergès et contributeurs. Découpage et classement : CoopMaths, Sébastien Lozano, Éric Elter et contributeurs. Adaptation de la présentation : Maths Explorer.
% CC BY-SA : https://coopmaths.fr/www/about/
% https://coopmaths.fr/alea/static/eam/2026/tex/eam_specifique_2026_06_asie_2_cor.tex
\begin{document}
\noindent\textit{APMEP / CoopMaths. CC BY-SA. Présentation adaptée par Maths Explorer.}\par\medskip
\subsection*{Partie A -- Gain d'expérience}
\begin{enumerate}
	\item Au niveau 0, le personnage a 20~points, il en gagne 80 lorsqu'il change de niveau, donc au niveau~1, le personnage a $20+80 = 100$ points d'expérience.

	\item Lorsque le personnage est au niveau 10, il a changé 10 fois de niveau, donc acquis $10\times80 = 800$ points d'expérience, qui s'ajoutent aux 20 du début de la partie.

	Arrivé au niveau 10, le personnage a donc plus de 800 points d'expérience, 820 pour être précis.

	\item Si le personnage a \np{2020} points d'expérience, cela signifie qu'il a \np{2000} points de plus qu'au début de la partie.

	Si on appelle $n$ le niveau atteint, il a changé $n$ fois de niveau, et on sait que \quad $n\times 80 = \np{2000}$, \quad donc $n = \dfrac{\np{2000}}{80} = \dfrac{200}{8} = \dfrac{100}{4} = 25$.

	Si le personnage a \np{2020} points d'expérience, il est donc au niveau 25.
\end{enumerate}

\subsection*{Partie B -- Gain d'énergie magique}
\begin{enumerate}
	\item Au niveau~0, le monstre neutralisé permet de récupérer 10~centilitres (cl) de potion magique, donc $v_0 = 10$.

	\item À chaque niveau, cette quantité augmente de 20~\% par rapport au niveau précédent.

	Augmenter une quantité de 20~\% revient à la multiplier par $\dfrac{100+20}{100} = \np{1.2}$.

	Donc, si $v_n$ est la quantité de potion magique récupérée au niveau $n$, \quad $v_{n+1} = \np{1.2} v_n$,\quad pour tout entier naturel $n$.

	\item $(v_n)$ est donc une suite géométrique de raison \np{1.2} et de premier terme $v_0 = 10$.

	\item L'expression explicite d'une suite géométrique $(v_n)$ de premier terme $v_0$ et de raison $q$ est :

	$\forall n \in \N$,\quad $v_n = v_0~q^n$,

	donc l'expression de $(v_n)$ en fonction de $n$ est $ v_n = 10 \times \np{1.2}^n$, pour tout entier naturel $n$.

	\item \begin{enumerate}
		\item On sait que \quad $v_5 = 10\times \np{1.2}^5 = \np{24,8832}$,

		donc $v_5\approx25$, arrondi à l'unité.

		\item Cela signifie que le monstre neutralisé à la fin du niveau 5 permettra au personnage de récupérer environ 25~cl de potion magique.
	\end{enumerate}

	\item $\np[~litre]{0.6} = \np[~cl]{60}$, donc on cherche pour quelle valeur de $n$ $v_n$ dépassera 60.

	Puisque $ v_n = 10 \times \np{1.2}^n$,\quad $v_n$ dépasse 60 si $\np{1.2}^n$ est plus grand que 6 : il faut donc attendre le niveau 10 pour que la quantité de potion magique dépasse \np{0.6}~litre.
\end{enumerate}

\needspace{10\baselineskip}

\end{document}
